\documentclass[10pt]{article}
\usepackage[margin=20mm]{geometry}
\usepackage{amsmath,amssymb,mathtools,fontspec,unicode-math,xurl,hyperref}
\IfFontExistsTF{DejaVu Serif}{\setmainfont{DejaVu Serif}}{\setmainfont{Latin Modern Roman}}
\setmathfont{latinmodern-math.otf}
\setlength{\parindent}{0pt}\setlength{\parskip}{6pt}\sloppy
\title{Scientific review of
Neck pinch singularities and Joyce conjectures in Lagrangian mean curvature flow with circle symmetry}
\author{M25-0288}
\date{27 October 2026}
\begin{document}
\maketitle
Jason D. Lotay and Gonçalo Oliveira
Journal of the European Mathematical Society; DOI: 10.4171/jems/1661
Reviewed source: arXiv:2305.05744v2, 20 December 2024, 27 pages.\par
\section*{Scope and assessment}
This report concerns the exact author manuscript identified above. The frozen main PDF has SHA256\par
`8f7340383e3e68edbe9ad993af7d938d1493df2fdcd71119413cec96116a482b`.\par
The ordinary publisher download returned a subscription page rather than a PDF. Equivalence with the publisher’s version of record has not been established. The manuscript describes itself as accepted in JEMS, with changes to the proofs of Propositions 3.4 and 4.2. Those changes are part of the version reviewed here.\par
Outcome: Major for P3; P1, P2 and P4 remain Minor. Other proof questions remain Needs Evidence. Four defects are established below: inconsistent preliminary normalizations, a false gradient evolution identity, an invalid first-contact argument in the convexity proof, and a missing stationary exception in Proposition 3.5. P3 concerns a proof step explicitly used to obtain compact finite-time singularities and the phase-order conclusion of a main theorem. A corrected convexity proof is supplied, with its assumptions verified separately; no counterexample to Theorems 1.1 or 1.4 is asserted. The ALE exterior argument and the passage from rescaled curves to the hypotheses of the imported singularity theorems require more precise justification; the present review neither certifies those steps nor asserts an unproved counterexample. The severity assessment concerns the affected derivations and their demonstrated dependencies. The scientific decision remains pending, and the human-review list is empty.\par
\section*{Coverage and mathematical structure}
All 27 PDF pages were read as extracted text. Actual page images were inspected for pages 2–4 and 8–27, including all six figures. The source TeX was additionally checked at the disputed normalization, evolution, convexity and area formulas. The archive returned by arXiv is a single gzip-compressed TeX file: the initial tar-extraction failure is preserved, and its subsequent correct decoding has a separate receipt.\par
The essential hypotheses are retained throughout: a complete ALE or ALF Gibbons–Hawking space, a finite set of centres, the potential \[\phi=m+\sum_{i=1}^k\frac1{2|x-p_i|},\qquad m\geq0,\] circle invariance, embeddedness and almost calibration where stipulated. A compact connected invariant Lagrangian is represented by an arc with centre endpoints; the noncompact examples have a different existence-at-infinity requirement. With Euclidean arclength \(s\), the reduction is \[\gamma_t=a\kappa N,\qquad a=\phi^{-1},\qquad \kappa=\theta_s,\qquad
 \kappa_t=(a\kappa)_{ss}+a\kappa^3.\] The direct moving-frame derivation on pp. 8–9 is consistent with this curvature equation. Also \[|\nabla\phi|\leq\sum_i(2r_i^2)^{-1}\leq2\phi^2,\qquad |\nabla a|\leq2;\] the signs and powers in the stated lower bound for \((a\kappa)_t\) follow.\par
The proof chain read was: preliminary convexity and disk variation (Section 3); localisation at centres, blow-up classification and Lawlor-neck identification (Section 4.1–4.2); finite-time barriers (Section 4.3–4.4); restart and the finite-centre argument (Section 5.1); convex-hull and stability alternatives (Section 5.2); long-time and current convergence (Section 5.3); and the distinction between chain order and phase order (Section 6). The perfect-Morse hypothesis in Theorem 1.4(d) and the allowance of equal limiting phases were not dropped. Figures 6.1–6.2 illustrate a nonconvex case and do not contradict the phase-order conclusion restricted by that hypothesis.\par
\section*{P1: preliminary Gibbons–Hawking normalizations}
Severity: Minor. Locations: p. 4, definition of \(\alpha_i\), (2.1), and Proposition 2.1(a),(c).\par
The manuscript defines \(|\xi|^2=\phi^{-1}\), \(\eta(\xi)=1\), with \(\ker\eta=\xi^\perp\), and \(\alpha_i=I_i(\phi^{-1}\eta)\), but writes \(g=\phi^{-1}\eta^2+\sum_i\alpha_i^2\). Choose a \(g\)-orthonormal coframe \(e^0,e^1,e^2,e^3\) adapted to \(\xi\). Then \(\eta=\sqrt{\phi}\,e^0\) and, up to irrelevant signs, \(\alpha_i=\phi^{-1/2}e^i\). Consequently the printed right-hand side is \[(e^0)^2+\phi^{-1}\bigl((e^1)^2+(e^2)^2+(e^3)^2\bigr),\] which differs from \(g\) whenever \(\phi\ne1\).\par
There is a second explicit inconsistency. Parts (a) and (c) together give \(g_E=\phi^{-2}\sum\alpha_i^2=\phi^{-1}\sum d\mu_i^2\). For the single-centre ALE potential \(\phi=1/(2r)\), this is \(2r\,\delta\), which is not flat. Indeed, for \(u=\frac12\log(2r)\), the three-dimensional conformal scalar-curvature calculation gives \[R_{e^{2u}\delta}
 =e^{-2u}\{-4\Delta_\delta u-2|\nabla u|^2\}
 =-\frac5{4r^3}.\] This supplies a local geometric test, rather than a reliance on a conventional choice of symbols.\par
A consistent orthonormal convention is \[\alpha_i=I_i(\phi^{-1/2}\eta),\qquad
 g_E=\phi^{-1}\sum_i\alpha_i^2,\qquad
 \mathrm{vol}_g=\phi^{-1/2}\eta\wedge\alpha_1\wedge\alpha_2\wedge\alpha_3.\] It yields the usual metric (2.4) and forms (2.5) already printed in the manuscript. Thus this is a repairable normalization defect in the preliminaries, not a counterexample to the later analysis formulated using (2.4). The positive factor in the printed volume form does not by itself change orientation; the issue is metric normalization. The source TeX confirms the powers, excluding an extraction artefact.\par
\section*{P2: the gradient evolution identity omits moving-metric terms}
Severity: Minor. Location: p. 8, the unnumbered identity following (3.2).\par
The asserted identity is \[(\partial_t+\Delta)|d\theta|^2=-2|\nabla d\theta|^2,
 \qquad \Delta=d^*d.\] A counterexample can be made within the graded, circle-invariant setting; a whole product torus, which is not graded, is unnecessary. In flat \(\mathbb C^2\), put \[R(t)^2=R_0^2-2t,\qquad
 L_t=\{(R(t)e^{iu},R(t)e^{iv}):u+v\pmod{2\pi}\in(-\delta,\delta)\},
 \quad 0<\delta<\pi/2 .\] This is a smooth open annulus, invariant under the tri-Hamiltonian action \((z_1,z_2)\mapsto(e^{i\alpha}z_1,e^{-i\alpha}z_2)\). It is a restriction of the normal mean curvature flow of the product of two shrinking circles. It has a global grading \(\theta=u+v+\pi\) on the chosen branch and angle oscillation \(2\delta<\pi\). Intrinsically, \[g_t=R(t)^2(du^2+dv^2),\quad
 |d\theta|^2=\frac2{R(t)^2},\quad
 \nabla d\theta=0,\quad
 \Delta|d\theta|^2=0,\quad
 \partial_t|d\theta|^2=\frac4{R(t)^4}>0.\] The displayed identity is therefore false even in flat ambient geometry and for a global grading. This is a local identity test on an open annulus, not a compact-flow counterexample to Theorem 1.4.\par
In Euclidean ambient space the missing contribution can be written \(2|A(\nabla\theta,\cdot)|^2\) in the standard moving-metric Bochner formula; curved ambient space also requires its curvature terms. The identity should be removed or replaced with a correctly stated formula. The subsequent curvature equation has a separate direct proof and is not invalidated by this example. No essential later use of the false equality was established in this review.\par
\section*{P3: the first-contact argument crosses its sign domain}
Severity: Major. Location: Proposition 3.4, pp. 10–11, (3.8)–(3.9).\par
Set \(f=a\kappa\). The manuscript obtains \(f_t\geq a f_{ss}-cf\) only on the set \(f\geq0\), then applies it at a first zero of \(f^\varepsilon=f+\varepsilon t\), with \(t_0>0\). At that point \[f(t_0,x_0)=-\varepsilon t_0<0,\] so the preceding sign-restricted inequality is unavailable. The text also replaces this value by zero. This invalidates the written contradiction, independently of the separate assertion about short-time positivity of the mean-curvature norm.\par
The statement has a direct repair under the smoothness and bounded-curvature assumptions invoked in this proof. Since \(a_t=(\partial_Na)f\), the exact curvature equation gives \[f_t=a f_{ss}+b f,\qquad b=\kappa\,\partial_Na+a\kappa^2.\] On the stipulated compact smooth time interval, \(|\nabla a|\leq2\), bounded \(\kappa\), and bounded \(a\) give \(|b|\leq B\). Take \(K>B\) and \(w=e^{-Kt}f\). Then \[w_t=a w_{ss}+(b-K)w.\] Initially \(w\geq0\), and its endpoint values are zero. If it had a negative value, a negative minimum over a compact preceding spacetime interval would lie at positive time and at an interior spatial point. There \(w_t\leq0\), \(w_{ss}\geq0\), and \((b-K)w>0\), a contradiction. Hence \(f\geq0\). The strong maximum principle on interior regions, where \(a>0\), gives strict interior positivity from the strictly positive initial curvature. Endpoint degeneracy does not interfere with the negative-minimum argument, since the endpoint values are zero.\par
This repairs the relevant convexity conclusion without assuming a maximum principle for \(|H|\). It applies on the smooth compact time intervals, with bounded curvature and zero endpoint values, stipulated in the source proof. It does not certify weaker endpoint regularity or the separate passage through singular times.\par
Dependency and impact. Proposition 3.4 is explicitly invoked in Theorem 4.13 (p. 19), which supplies Theorem 4.14 and the compact examples in Theorem 1.1(b). It is also explicitly invoked in Corollary 5.5 (p. 22); that convex-hull result is then used for the phase-order conclusion of Theorem 5.9(b) (p. 24), and hence Theorem 1.4(d) (p. 3). The written first-contact contradiction is therefore a demonstrated invalid derivation tied directly to central contributions. This supports Major for P3. The replacement proof above restores this convexity step under its stated assumptions; repair verification is separate from the severity of the original argument, and does not certify the remaining global proof.\par
\section*{P4: Proposition 3.5 includes stationary straight lines}
Severity: Minor. Location: p. 11, Proposition 3.5.\par
The proposition assumes only \(\kappa_0\geq0\), yet concludes \(\kappa_t>0\) for every \(t\in[0,T]\). Take a single centre at the origin and \[\phi(x)=\frac1{2|x|},\qquad \gamma_t(s)=(s,1,0),\quad s\in\mathbb R.\] This embedded planar curve avoids the centre, has \(\kappa_t=0\), and satisfies \(\partial_t\gamma=a\kappa N=0\) for all time. It meets every hypothesis of this auxiliary proposition and violates its conclusion even for \(t>0\). Inclusion of \(t=0\) is an additional immediate problem for initially nonnegative curvature with zeros.\par
A strict-positivity statement must exclude the identically straight solution and start at \(t>0\); in a noncompact formulation it must also specify the conditions allowing the relevant global maximum principle. These are necessary corrections, not a certification of every possible noncompact extension. The manuscript introduces this proposition for possible future study, and no dependence of the principal compact theorem on its overbroad form was identified.\par
\section*{Area convention and the finite-time barrier}
Definition 4.9 calls the pacman disk a planar region, whereas Proposition 4.10 uses an area derivative derived from the Kähler area of a lift. These agree only after specifying the weight. For a planar region \(D_t\), put \[E(t)=\int_{D_t}dA_E,\qquad B(t)=\int_{D_t}\phi\,dA_E .\] With consistent orientation and only \(\gamma_t\) moving, \[E'(t)=\int_{\gamma_t}a\kappa\,ds,\qquad
 B'(t)=\int_{\gamma_t}\kappa\,ds
       =\theta_t(p_+)-\theta_t(p_-).\] The latter is also the Kähler area density supplied by (2.5) on the plane \(\mu_1=\mathrm{constant}\). The centre singularities \(1/r\) are locally integrable in two dimensions.\par
Here is an exact diagnostic for the Euclidean interpretation. Let the three centres in the plane be \(p_-=(-1,0)\), \(p_+=(1,0)\), \(p=(0,1/2)\), take \(m=10\), and let \(\gamma_0(x)=(x,1-x^2)\), \(-1\leq x\leq1\). The lower two straight segments run from \(p_-\) to \(p\) to \(p_+\). Their graph is \(\frac12(1-|x|)\); the upper parabola lies strictly above it between the endpoints. The arc is almost calibrated since its tangent-angle oscillation is \(2\arctan2<\pi\). For any smooth flow from this datum the initial Euclidean variation is \[E'(0)=-\int_{-1}^1\frac{2\,dx}{\phi(x,1-x^2)(1+4x^2)}
 \geq-\frac{2\arctan2}{10}>-2\arctan(1/2)
 =\theta_+-\theta_- .\] Thus (4.10) cannot use Euclidean area. In contrast \(B'(0)=-2\arctan2\leq-2\arctan(1/2)\), as required. The strict comparison follows, without numerics, from \(\arctan2<2\) and \(\arctan(1/2)\geq(1/2)/(1+1/4)=2/5\).\par
The natural weighted-area interpretation repairs this specific issue. It is recorded as a clarification with a conditional counterexample, not as a fifth unconditional theorem defect. For general triads made of broken straight segments, the proof’s assertion that the chord lines \(\ell_\pm\) are parts of the barrier also needs explanation; the convex-polygon application has extra angle-order information. No verified counterexample to that full application is asserted.\par
\section*{Unresolved dependencies and rejected overstatements}
\subsection*{ALE exterior control}
On p. 15, the proof of Proposition 4.3 infers uniform equivalence with ordinary curve shortening outside a compact set from finiteness of the centre set. For ALE spaces, \[\phi(x)\sim\frac{k}{2|x|},\qquad a(x)\sim\frac{2|x|}{k}\longrightarrow\infty .\] An upper uniform comparison of coefficients is false. For ALF with \(m>0\), \(a\to1/m\), so the two cases differ. A weighted parabolic existence and asymptotic-preservation argument, or a precise match to the hypotheses of the cited work of Su, is needed for the noncompact branch. The full hypotheses of that external existence theorem were not accessible through the checked normal links. This is Needs Evidence concerning a core proof step, not a proven failure of existence.\par
\subsection*{Blow-up and imported theorems}
Theorem 1.1 of Lambert–Lotay–Schulze, arXiv:1901.05383v2, requires an ancient exact almost calibrated solution with uniformly bounded area ratios and angle, and a blow-down consisting of two transverse multiplicity-one planes. Its conclusion classifies that ancient solution. Its statement and the adjoining discussion of exactness of Type II blow-ups were checked on pp. 1–2. A graphical projected curve alone should not be substituted for this entire list; the transversality, multiplicity and actual blow-down must be supplied.\par
The Hopf model has \(\pi_H(\rho z)=\rho^2\pi_H(z)\). A base dilation \(\lambda\) therefore corresponds to ambient dilation \(\sqrt\lambda\) and the usual MCF time scale \(\lambda^{-1}\), not \(\lambda^{-2}\). In the one-centre equation \(a=2r\), the base-space/time rescaling used in (4.8) has an extra factor \(\lambda^{-1}\) in its velocity equation. It can be useful for a stationary spatial limit, but is not itself a standard parabolic tangent-flow scaling. The record-minimum condition (4.7), imposed after an arbitrary sequence of times is specified, also needs a selection or interpolation argument. These observations identify what the bridge to a tangent flow must prove; they do not show that the Lawlor-neck conclusion is false.\par
Theorems 1.1–1.2 of Lotay–Schulze–Székelyhidi, arXiv:2208.11054v1, were checked on pp. 1–2. They give tangent-flow uniqueness and a restart from a \(C^1\) immersed limit when the stipulated zero-Maslov/rational and two-plane tangent hypotheses hold at the first singular time. Hence an objection that the citation contains no restart theorem is rejected. The current manuscript’s localisation from a compact Calabi–Yau ambient space to the bounded Gibbons–Hawking region still needs its analytic justification. Neither external paper was independently reproved in full.\par
\subsection*{Currents and convergence}
Uniform convergence of sets alone would not imply current convergence with multiplicity. Here the almost-calibrated graphical description supplies additional length and multiplicity control, so that generic objection is not a counterexample. Likewise, the \(O(\varepsilon)\) swept-volume estimate in Proposition 5.11 is stronger than needed. For a compact smooth MCF segment, the area identity and Cauchy–Schwarz give \[\int_{T-\varepsilon}^T\!\int_{L_t}|H|
 \leq\bigl(\varepsilon\,\mathrm{Area}(L_0)\bigr)^{1/2}
       \left(\int_{T-\varepsilon}^T\!\int_{L_t}|H|^2\right)^{1/2}
 \longrightarrow0.\] This supports the required continuity once the limiting current and restart have been justified. The complete continuation and convergence package is not independently certified by this bounded repair.\par
\section*{Supporting checks}
The primary executor performed the scientific reasoning, exact calculations, current PDF parsing, compilation and actual image inspection. The independent Python/SymPy script checked the scalar-curvature computation, the graded-annulus evolution, the exact \(f\)-equation, the stationary line and the area/scaling diagnostics. Its actual execution exited with code 0 and all assertions passed. It is reviewer-authored; no author code was run and no PDE simulation was used.\par
The report is built with the previously approved Tectonic executable in an isolated staging cache, with a warm-up followed by an offline `–only-cached –untrusted` final build. Exact TeX, compiler, dependency and PDF hashes, real commands, exit codes and logs are bound in the build receipts. Every report page is parsed, rendered and actually inspected before completion is recorded; rendering-only false records and the source-archive extraction failure are preserved. The source-page record contains only actually viewed images. The two external PDFs have separate acquisition, parse, render and actual-view receipts.\par
The native assistant `/root/binding\_checklist` supplied only a non-scientific contract-field checklist. No scientific reasoning, proof, computation, compilation, parsing or page viewing was delegated. Recovered historical auxiliary source-search roles remain in the frozen input evidence. Controller-supplied UI observations are Chat pressed, Work not pressed, official GPT-6 checked, Pro 5 of 5, and bigbroliu/Personal checked. Runtime self-identification is GPT-6 Astra Pro; the backend was not independently verified. The initial recovery did not perform independent acceptance. A subsequent focused independent acceptance of this frozen paper verified P1–P4, corrected the severity and dependency statement for P3, and retained the exact-version limitations; its current bindings and actual execution facts are recorded in the existing verification files.\par
The unavailable publisher PDF and separately released supplements/code are material-access limitations, not mathematical defects. All statements of this report refer to the pinned version. Imported results and unresolved steps above remain explicitly bounded.\par
\end{document}
